\chapter{Hilbert Spaces}

\section{Orthogonality and direct sums}

\begin{noteproposition}[Orthogonal density criterion]
Let $H$ be a Hilbert space and let $M\subseteq H$ be a linear subspace.  Then
\[
  \overline M=H
  \quad\Longleftrightarrow\quad
  M^\perp=\{0\}.
\]
\end{noteproposition}

\begin{proof}
If $\overline M=H$, continuity of the inner product makes every vector
orthogonal to $M$ orthogonal to all of $H$, hence zero.  Conversely, if
$\overline M\ne H$, the projection theorem gives a nonzero component of some
$x\notin\overline M$ in $(\overline M)^\perp=M^\perp$.
\end{proof}

\begin{noteproposition}
For any nonempty subset $S$ of an inner-product space,
\[
  S^\perp=\{x:\inner{x}{s}=0\text{ for every }s\in S\}
\]
is a closed linear subspace.  In a Hilbert space,
\[
  S^\perp=(\overline{\Span S})^\perp,
  \qquad
  S^{\perp\perp}=\overline{\Span S}.
\]
Moreover, $S\cap S^\perp=\{0\}$ whenever $S$ is a subspace.
\end{noteproposition}
\begin{notetheorem}[Orthogonal direct-sum theorem]
Let $M$ be a closed linear subspace of a Hilbert space $H$.  Then
\[
  H=M\oplus M^\perp.
\]
Thus every $x\in H$ has a unique decomposition
\[
  x=P_Mx+(x-P_Mx),
  \qquad
  P_Mx\in M,
  \quad x-P_Mx\in M^\perp.
\]
\end{notetheorem}

\section{Inner products and projections}

\begin{notedefinition}
An inner product on a real vector space is a symmetric positive-definite
bilinear form.  On a complex vector space it is a positive-definite
sesquilinear form, conjugate linear in one variable and linear in the other.
The induced norm is
\[
  \|x\|=\sqrt{\inner{x}{x}}.
\]
A complete inner-product space is a \emph{Hilbert space}.
\end{notedefinition}

\begin{noteproposition}
Every inner product satisfies the Cauchy--Schwarz inequality
\[
  |\inner{x}{y}|\le\|x\|\,\|y\|
\]
and the parallelogram identity
\[
  \|x+y\|^2+\|x-y\|^2
  =2\|x\|^2+2\|y\|^2.
\]
Conversely, a norm comes from an inner product if and only if it satisfies the
parallelogram identity; the inner product is recovered by polarization.
\end{noteproposition}
\begin{notetheorem}[Projection theorem]
Let $K$ be a nonempty closed convex subset of a Hilbert space $H$.  For every
$f\in H$ there is a unique $u\in K$ satisfying
\[
  \|f-u\|=\dist(f,K).
\]
It is characterized by the variational inequality
\[
  \operatorname{Re}\inner{f-u}{v-u}\le0
  \qquad(v\in K).
\]
We write $u=P_Kf$.
\end{notetheorem}

\begin{proof}
Let $d=\dist(f,K)$ and choose a minimizing sequence $(u_n)\subseteq K$.
Convexity gives $(u_n+u_m)/2\in K$.  The parallelogram identity yields
\[
 \|u_n-u_m\|^2
 =2\|f-u_n\|^2+2\|f-u_m\|^2
  -4\left\|f-\frac{u_n+u_m}{2}\right\|^2,
\]
which tends to zero.  Thus $(u_n)$ is Cauchy and converges to some $u\in K$;
continuity gives $\|f-u\|=d$.

For $v\in K$ and $0<t\le1$, the point $u+t(v-u)$ lies in $K$.  Minimality
implies
\[
  \|f-u\|^2
  \le\|f-u-t(v-u)\|^2
  =\|f-u\|^2-2t\operatorname{Re}\inner{f-u}{v-u}
    +t^2\|v-u\|^2.
\]
Divide by $t$ and let $t\downarrow0$ to obtain the inequality.  Conversely,
the inequality gives, for $v\in K$,
\[
  \|f-v\|^2
  =\|f-u\|^2+\|u-v\|^2
    +2\operatorname{Re}\inner{f-u}{u-v}
  \ge\|f-u\|^2.
\]
If $u_1,u_2$ both satisfy the variational inequalities, use $v=u_2$ in the
first and $v=u_1$ in the second; adding gives
$\|u_1-u_2\|^2\le0$.
\end{proof}
\begin{noteproposition}[Nonexpansiveness of metric projection]
For every nonempty closed convex $K\subseteq H$,
\[
  \|P_Kf-P_Kg\|\le\|f-g\|
  \qquad(f,g\in H).
\]
\end{noteproposition}

\begin{proof}
Apply the variational inequality to $P_Kf$ with competitor $P_Kg$ and to
$P_Kg$ with competitor $P_Kf$.  Adding gives
\[
  \|P_Kf-P_Kg\|^2
  \le\operatorname{Re}\inner{P_Kf-P_Kg}{f-g}
  \le\|P_Kf-P_Kg\|\,\|f-g\|.
\]
\end{proof}

\begin{notecorollary}[Orthogonal projections]
If $M$ is a closed linear subspace, then $P_M$ is linear, bounded, and
$\|P_M\|\le1$.  It is characterized by
\[
  P_Mf\in M,
  \qquad
  f-P_Mf\in M^\perp.
\]
If $M\ne\{0\}$, then $\|P_M\|=1$.
\end{notecorollary}

\begin{proof}
The variational inequality may be tested with $u\pm v$ and, in the complex
case, $u\pm iv$ for $v\in M$, so it reduces to orthogonality.  Uniqueness of
the orthogonal decomposition gives linearity.  Nonexpansiveness gives the norm
bound.
\end{proof}
\section{Riesz representation}

\begin{notetheorem}[Riesz representation theorem]
For every $\ell\in H^*$ there is a unique $y\in H$ such that
\[
  \ell(x)=\inner{x}{y}
  \qquad(x\in H)
\]
(with the variables adjusted to the chosen complex convention), and
$\|\ell\|=\|y\|$.
\end{notetheorem}

\begin{proof}
If $\ell=0$, take $y=0$.  Otherwise $M=\Ker\ell$ is a closed proper subspace.
Choose $0\ne z\in M^\perp$.  For each $x\in H$, the vector
\[
  x-\frac{\ell(x)}{\ell(z)}z
\]
belongs to $M$, hence is orthogonal to $z$.  Therefore
\[
  \ell(x)=\frac{\ell(z)}{\|z\|^2}\inner{x}{z},
\]
which has the required form.  Cauchy--Schwarz gives
$\|\ell\|\le\|y\|$, while testing at $x=y/\|y\|$ gives equality.
\end{proof}
\section{Orthogonal sums and orthonormal bases}

\begin{notedefinition}
Let $(E_n)_{n\ge1}$ be closed subspaces of a Hilbert space $H$.  Their
Hilbert direct sum is
\[
  \bigoplus_{n=1}^{\infty}E_n
  =\left\{(u_n):u_n\in E_n,
                 \ \sum_{n=1}^{\infty}\|u_n\|^2<\infty\right\}
\]
with inner product
$\inner{(u_n)}{(v_n)}=\sum_n\inner{u_n}{v_n}$.
\end{notedefinition}

\begin{notelemma}[Orthogonal series]
If $(v_n)$ is an orthogonal sequence in $H$ and
$\sum_n\|v_n\|^2<\infty$, then $\sum_nv_n$ converges in $H$ and
\[
  \left\|\sum_{n=1}^{\infty}v_n\right\|^2
  =\sum_{n=1}^{\infty}\|v_n\|^2.
\]
\end{notelemma}

\begin{proof}
For $m>n$, orthogonality gives
\[
  \left\|\sum_{k=n+1}^m v_k\right\|^2
  =\sum_{k=n+1}^m\|v_k\|^2,
\]
so the partial sums are Cauchy.  The norm identity follows by passage to the
limit.
\end{proof}

\begin{notetheorem}[Projections onto mutually orthogonal subspaces]
Let $(E_n)$ be mutually orthogonal closed subspaces and let $P_n=P_{E_n}$.  For
all $u\in H$,
\[
  \sum_{n=1}^{\infty}\|P_nu\|^2\le\|u\|^2,
\]
and the series $\sum_nP_nu$ converges.  If
$E=\overline{\Span\bigcup_nE_n}$, then
\[
  u=P_Eu=\sum_{n=1}^{\infty}P_nu,
  \qquad
  \|P_Eu\|^2=\sum_{n=1}^{\infty}\|P_nu\|^2.
\]
\end{notetheorem}

\begin{proof}
Finite orthogonality gives
\[
 \left\|\sum_{n=1}^N P_nu\right\|^2
 =\sum_{n=1}^N\|P_nu\|^2
 =\inner{u}{\sum_{n=1}^N P_nu}
 \le\|u\|\left\|\sum_{n=1}^N P_nu\right\|.
\]
This yields Bessel's bound and convergence.  The limit lies in $E$, and the
difference from $u$ is orthogonal to every $E_n$, hence to $E$; it is
therefore $P_Eu$.
\end{proof}
\begin{notedefinition}
A sequence $(e_n)$ is an \emph{orthonormal basis} of $H$ if
\[
  \inner{e_n}{e_m}=\delta_{nm}
  \qquad\text{and}\qquad
  \overline{\Span\{e_n:n\ge1\}}=H.
\]
\end{notedefinition}

\begin{notetheorem}[Fourier expansion and Parseval identity]
If $(e_n)$ is an orthonormal basis, then for every $u\in H$,
\[
  u=\sum_{n=1}^{\infty}\inner{u}{e_n}e_n,
  \qquad
  \|u\|^2=\sum_{n=1}^{\infty}|\inner{u}{e_n}|^2.
\]
Conversely, for every $(a_n)\in\ell^2$ the series $\sum_na_ne_n$ converges to
a unique $u\in H$ satisfying
\[
  \inner{u}{e_n}=a_n,
  \qquad
  \|u\|^2=\sum_n|a_n|^2.
\]
\end{notetheorem}

\begin{notetheorem}
Every separable Hilbert space has a countable orthonormal basis.
\end{notetheorem}

\begin{proof}
Choose a countable dense set, discard vectors already in the span of their
predecessors, and apply the Gram--Schmidt process.  The resulting orthonormal
sequence has dense linear span because it spans the same finite-dimensional
spaces at each stage as the retained dense sequence.
\end{proof}
